\section{Uso de IA}
Considero que uno de los mayores provechos que le pude sacar a la IA fue para mejorar un aspecto de los streams. Al principio, había pensado los streams como solamente dos, uno para cálculo y otro para copias.
Revisando la clase teórica de CUDA Streams, me di cuenta que se podría lograr una mayor optimización haciendo un vector de streams homogéneos, es decir, que hagan las mismas operaciones y no se diferencien entre, por ejemplo, cálculos y copias.

\begin{lstlisting}
cudaStream_t stream_calculo, stream_copia;
HANDLE_ERROR(cudaStreamCreate(&stream_calculo));
HANDLE_ERROR(cudaStreamCreate(&stream_copia));
...
HANDLE_ERROR(cudaStreamWaitEvent(stream_copia, ev_fin_kernel[actual], 0));
HANDLE_ERROR(cudaEventRecord(ev_inicio_copia, stream_copia));
HANDLE_ERROR(cudaMemcpyAsync(h_grilla[siguiente], d_grilla[siguiente],
                           bytes, cudaMemcpyDeviceToHost,
                           stream_copia));
HANDLE_ERROR(cudaEventRecord(ev_fin_copia, stream_copia));

HANDLE_ERROR(cudaEventRecord(ev_inicio_kernel[siguiente], stream_calculo));
difusion<<<grid, bloque, bytes_compartidos, stream_calculo>>>(
  d_grilla[siguiente], d_grilla[actual], N, distancia);
HANDLE_ERROR(cudaPeekAtLastError());
HANDLE_ERROR(cudaEventRecord(ev_fin_kernel[siguiente], stream_calculo));
\end{lstlisting}

Para mejorarlo y hacerlo como la cátedra, le pedí a Claude Opus 5.5 y a Claude Sonnet 5.5 lo siguiente:

\begin{block}{Prompt}
\textit{Te adjunto un programa CUDA que simula la difusión de calor en una grilla N×N y guarda cada paso de la simulación en un archivo binario. Usa dos streams, uno para el kernel y otro para la copia D2H, coordinados con cudaStreamWaitEvent, con memoria pinned y doble buffer, de modo que la escritura a disco de un paso se superpone con el cálculo y la copia del siguiente. Modificalo para que siga el esquema de la filmina que adjunto, que sería un vector de N (usá 5 para probar) streams, donde cada stream procesa su porción de la grilla, de modo que la copia de una porción se superponga con el cálculo de otra. La escritura a disco se tiene que seguir superponiendo con el trabajo de la GPU. Los archivos generados tienen que quedar idénticos byte a byte a los del programa original. Dejá los mismos tres parámetros y que imprima los tiempos con el mismo formato. Explicá las decisiones que tomaste.}
\end{block}

\begin{table}[H]
  \centering
  \small
  \begin{tabular}{p{4cm}p{4cm}p{4cm}p{4cm}}
    \toprule
    Criterio & Dos streams (mi versión) & Opus & Sonnet \\
    \midrule
    Pensó durante & -- & 8' 30" & 11' \\
    Compila & Si & Si & Si \\
    Frames idénticos a la versión base (N = 1024) & -- & Sí & Sí \\
    Streams & 2 (cálculo y copia) & 5 & 5 \\
    División de la grilla & No divide & Bandas de filas de bloques completas & Bandas de filas \\
    Cambio en el kernel & -- & Recibe la primera fila de bloques de la banda & Recibe la fila inicial y final, y cambia la carga del halo inferior \\
    Sincronización & La copia espera al kernel con cudaStream\-WaitEvent & Cadena de eventos entre todos los kernels & Eventos con las dos bandas vecinas \\
    Trabajo de la GPU mientras se escribe un paso & Copia del siguiente y cálculo del subsiguiente & Copia del siguiente y cálculo del subsiguiente & Copia del siguiente y cálculo del subsiguiente \\
    Líneas de código (sin comentarios) & 234 & 291 & 279 \\
    Prompts hasta que funcionó & -- & 1 & 1 \\
    \bottomrule
  \end{tabular}
  \caption{Comparación de la versión con dos streams y las soluciones con un vector de streams.}
  \label{tab:ia_comparacion}
\end{table}

Las dos soluciones generan frames idénticos a los de la versión base y detectaron por su cuenta la dependencia del halo entre bandas, pero la resuelven distinto. Opus alinea las bandas a filas de bloques completas, así el kernel solo necesita saber dónde empieza su banda, y encadena todos los kernels con eventos, por lo que nunca se ejecutan dos a la vez. Sonnet sincroniza cada banda solo con sus dos vecinas, pero modificó más el kernel para cargar el halo inferior, y uno de sus comentarios justifica una dependencia con un argumento incorrecto, aunque sin afectar el resultado.

Opus además agrupa los datos de las bandas en una estructura con funciones propias, lo que acorta el bucle principal a cambio de más código.

La versión que entrego es más liviana que las dos y se acerca más a la de Sonnet. No usa estructuras ni funciones auxiliares y se le pasa al kernel la fila de inicio y de fin de cada banda. A diferencia de ambas, sincroniza con un único cudaDeviceSynchronize por paso en lugar de eventos entre streams, que resuelve el halo de forma más simple y no cuesta tiempo, porque la GPU termina antes que la escritura.

\begin{table}[H]
  \centering
  \begin{tabular}{lrrrrr}
    \toprule
    Versión & Total (s) & Escritura (s) & Espera (s) & Kernel (ms/paso) & D2H (s) \\
    \midrule
    Dos streams & 1.5465 & 1.5455 & 0.0001 & 0.0317 & 0.0043 \\
    Opus        & 1.5647 & 1.5626 & 0.0001 & 0.0606 & 0.0052 \\
    Sonnet      & 1.5854 & 1.5811 & 0.0001 & 0.0639 & 0.0053 \\
    \bottomrule
  \end{tabular}
  \caption{Tiempos de las tres versiones en una misma ejecución del script de comparación (N = 1024, 30 pasos, 256 hilos por bloque). Opus y Sonnet con 5 streams.}
  \label{tab:ia_tiempos}
\end{table}

Fuera de la escritura, la versión con dos streams suma 1.0 ms, contra 2.1 ms de Opus y 4.3 ms de Sonnet, lo que es algo esperable ya que en cada paso se encolan cinco kernels y cinco copias en lugar de uno y una. Al ser una sola ejecución, estas diferencias de pocos milisegundos pueden variar entre corridas. Con el vector de streams el kernel tarda alrededor del doble por paso y la copia un 20\% más, posiblemente por lanzar cinco kernels y cinco copias más chicas en cada paso, pero como el kernel ocupa menos del 0.1\% del tiempo total no se nota.

\begin{table}[H]
  \centering
  \small
  \begin{tabular}{p{1.5cm}p{5.5cm}p{1.4cm}p{6cm}}
    \toprule
    Asistente & Prompt & Exitosa & Razones y solución \\
    \midrule
    \multicolumn{4}{l}{\textit{Uso para los requerimientos del enunciado:}} \\
    \midrule
    Claude & Generar scripts para un barrido de parámetros en el cluster hexa. & Sí & \\
    Claude & Calcular el speedup entre el secuencial y la mejor configuración CUDA. & Parcial & Usó el tiempo total sin aclarar que la teoría toma solo el kernel y las transferencias. Mantuve el tiempo total como criterio principal, justificando por qué, e informé también el speedup con el criterio de la teoría. \\
    Claude & Analizar los resultados e identificar aspectos ocultos. & Parcial & Detectó que el secuencial y la versión CUDA se medían distinto. Agregué la medición dentro del secuencial y repetí los barridos. \\
    Claude & Probar en el cluster cada paso de la versión con streams. & Sí & Encontró un error de compilación por el nombre de un evento y que el tiempo del kernel se sumaba dos veces. Corregí los dos. \\
    Claude & Script que compare los frames de la versión con streams contra la versión sin streams. & Sí & \\
    \midrule
    \multicolumn{4}{l}{\textit{Uso para la comparación de soluciones:}} \\
    \midrule
    Claude Opus 5.5 & Llevar la versión de dos streams a un vector de streams según la teoría (conversación nueva). & Sí & Frames idénticos. \\
    Claude Sonnet 5.5 & Llevar la versión de dos streams a un vector de streams según la teoría (conversación nueva). & Sí & Frames idénticos. \\
    \bottomrule
  \end{tabular}
  \caption{Interacción con los motores de IA.}
  \label{tab:ia}
\end{table}

Reconozco, de todas formas, que en este problema el vector de streams no aporta una mejora sobre la versión con dos streams. La ganancia frente a la versión sin streams viene de superponer la escritura a disco con el trabajo de la GPU, algo que los dos streams ya lograban. Dividir la grilla en bandas no agrega más: no hay copias H2D en cada paso, la GPU ya quedaba oculta detrás de la escritura, y lo máximo que podría ahorrar superponer la copia de una banda con el cálculo de otra es el tiempo del kernel, unos 0.03 ms por paso. A cambio, cada paso encola cinco kernels y cinco copias en lugar de uno y una, un costo del mismo orden que esa ganancia. El esquema de la teoría conviene cuando las transferencias y el cálculo pesan parecido y están en el camino crítico, que no es el caso cuando el límite es la escritura en la CPU.